% Section 1: Background and Motivation
\section{Background and Motivation}

\subsection{Video Super-Resolution for Long Videos}

Video super-resolution (VSR) has advanced rapidly with the adoption of diffusion-based and transformer-based
architectures. Where earlier CNN-based methods upscaled fixed short clips (typically 7--32 frames), modern
methods such as MGLD-VSR \cite{he_mgld_2024} and Upscale-A-Video \cite{zhou_upscale_2024} can process
substantially longer sequences in latent or sliding-window fashion, enabling deployment on real-world footage
that runs into minutes rather than seconds.

This shift toward long-video SR introduces an evaluation gap. The standard VSR benchmarks --- REDS (100-frame
clips at $720\!\times\!1280$), Vimeo-90K (7-frame septuplets), Vid4, UDM10, and SPMCS --- were designed
for short-clip evaluation and report full-reference metrics (PSNR, SSIM, LPIPS) against paired HR ground
truth. When applied to long videos in the wild, three fundamental limitations emerge:

\begin{itemize}
    \item \textbf{No HR ground truth.} Real deployment videos have no paired high-resolution reference,
          making full-reference metrics (PSNR, SSIM, LPIPS, DISTS) inapplicable.
    \item \textbf{Adjacent-frame temporal metrics miss long-range artefacts.} Existing temporal metrics
          --- tOF and tLP from TecoGAN \cite{chu_tecogan_2020}, and E*warp from DOVE \cite{yang_dove_2024}
          --- are computed at a single adjacent-frame gap ($k{=}1$). A method that is locally smooth
          but drifts globally over minutes is undetectable by these measures.
    \item \textbf{Smoother-output bias in learned representations.} DINOv2-based consistency metrics and
          LPIPS favour methods whose outputs change less in feature space between frames. Diffusion-based
          methods produce detail-rich outputs that score worse on these metrics despite being visually
          superior, a structural bias that misrepresents quality for modern SR architectures.
\end{itemize}

This thesis addresses the evaluation gap: long-video SR (videos longer than one minute) requires
no-reference metrics that capture temporal consistency across multiple time scales and are robust to
the smoother-output bias of learned representations.

\subsection{Existing Metric Landscape}

The existing metric ecosystem can be organised into five families:

\begin{enumerate}
    \item \textbf{Full-reference pixel/perceptual} (PSNR, SSIM, LPIPS, DISTS) --- require HR ground truth;
          inapplicable to deployment-time long-video evaluation.
    \item \textbf{No-reference single-aspect} (CLIP-IQA \cite{wang_clipiqa_2023}, MUSIQ \cite{ke_musiq_2021},
          NIQE, BRISQUE) --- measure per-frame quality; no temporal dimension.
    \item \textbf{Temporal pixel-level} (tOF, tLP from TecoGAN; E*warp from DOVE) --- adjacent-frame
          ($k{=}1$) optical-flow consistency; systematically miss long-range temporal structure.
    \item \textbf{Semantic video quality} (VBench \cite{huang_vbench_2024}, VBench-2.0
          \cite{huang_vbench2_2025}) --- multi-aspect benchmarks covering subject consistency,
          Human\_Identity, Human\_Anatomy, and others. These benchmarks were designed for
          text-to-video generation evaluation; thirteen of VBench-2.0's eighteen dimensions
          require text prompts and are therefore inapplicable to super-resolution, while the
          remaining prompt-independent dimensions operate at clip level and average per-clip
          scores up to a single video number, removing the cross-clip drift signal that
          long-video consistency assessment requires. Content-dependent failure modes are
          documented in Section~\ref{sec:preliminary}.
    \item \textbf{Disentangled aesthetic-technical NR} (DOVER \cite{wu_dover_2023}) ---
          a learned video quality predictor whose final score is a combination of two
          disentangled convolutional branches: an aesthetic branch trained on perception-driven
          labels and a technical branch trained on distortion-driven labels. DOVER is per-video
          and produces a single aggregate; it does not specifically model cross-frame
          temporal drift.
\end{enumerate}

The thesis problem is that no existing metric satisfies all four desiderata simultaneously:
(1)~no-reference, (2)~multi-time-scale temporal coverage, (3)~robust to smoother-output bias from
learned representations, and (4)~reliable across diverse long-video content including close-up scenes.

% Section 2: Literature Review
\section{Literature Review}
\label{sec:litreview}

\subsection{Video Super-Resolution Methods}

The two primary VSR methods evaluated throughout this thesis proposal represent the two
dominant stylistic families in modern super-resolution:

\begin{itemize}
    \item \textbf{MGLD-VSR} \cite{he_mgld_2024} (ECCV 2024) applies masked guided latent
          diffusion to progressively refine super-resolved frames. It is characterised by
          sharp, detail-preserving output and exact reproduction of published DOVE benchmark
          numbers on UDM10 (PSNR 24.23, SSIM and LPIPS to four decimal places), providing
          a reliable evaluation baseline.
    \item \textbf{Upscale-A-Video (UAV)} \cite{zhou_upscale_2024} (CVPR 2024) uses a
          text-conditioned temporal diffusion model inspired by text-to-video generation.
          It produces smoother, perceptually cleaner output at the cost of some
          high-frequency detail. UAV reproduces DOVE UDM10 results within a consistent
          $+1.33$~dB PSNR offset attributable to a minor degradation pipeline
          configuration difference, confirming the evaluation setup is correct.
\end{itemize}

The associated evaluation framework is provided by \textbf{DOVE} \cite{yang_dove_2024}
(NeurIPS 2025), the current VSR benchmark suite. DOVE combines the UDM10 and SPMCS test sets
with a multi-metric evaluation framework that includes E*warp for temporal consistency. DOVE
does not address long-video evaluation or multi-scale temporal metrics, and is therefore used
in this proposal as a baseline framework rather than as a long-video evaluation target.

\subsection{Temporal Consistency Metrics}

Three families of temporal consistency metrics have been used in the video super-resolution
and video generation literature. Each was originally designed for a different evaluation
target; together they form the temporal-evaluation toolkit currently available to a
super-resolution practitioner.

\begin{itemize}
    \item \textbf{tOF and tLP} (TecoGAN, Chu et al.\ \cite{chu_tecogan_2020}) measure
          per-frame optical-flow consistency (tOF) and LPIPS-based perceptual warping error
          (tLP). Both are defined at a single adjacent-frame gap ($k{=}1$). The preliminary
          work in Section~\ref{sec:preliminary} demonstrates that tOF ordering inverts
          between $k{=}1$ and $k{\geq}10$: the smoother method (UAV) wins at $k{=}1$ but
          loses to the detail-preserving method (MGLD-VSR) at long range.
    \item \textbf{E*warp} \cite{yang_dove_2024} extends warp-based error evaluation as part
          of the DOVE benchmark; it remains a $k{=}1$-dominant metric and shares the
          adjacent-frame limitation of tOF.
    \item \textbf{VBench} \cite{huang_vbench_2024} and \textbf{VBench-2.0}
          \cite{huang_vbench2_2025} provide multi-aspect evaluation of generated video,
          including subject and background consistency (DINOv2 cosine similarity between
          adjacent frames), Human Identity (slow-fast ArcFace clip adapter), and Human
          Anatomy (anomaly ViT). The preliminary work documents applicability gaps for
          super-resolution evaluation: DINOv2-based consistency metrics exhibit
          smoother-output bias and the Human Anatomy dimension has content-dependent
          calibration failure on close-up footage.
\end{itemize}

\subsection{No-Reference Perceptual Metrics}

Three no-reference perceptual metrics are commonly applied to super-resolution outputs when
high-resolution ground truth is unavailable. Each is single-frame or per-clip in scope and
none specifically targets long-range cross-frame consistency, but together they form the
no-reference perceptual evaluation toolkit relevant to deployment-time super-resolution
assessment.

\begin{itemize}
    \item \textbf{CLIP-IQA} \cite{wang_clipiqa_2023} evaluates per-frame image quality using
          CLIP embeddings trained on diverse natural content. Unlike LPIPS or DINOv2,
          CLIP-IQA does not systematically penalise high-frequency detail in
          diffusion-style super-resolution outputs, making it a suitable component of a
          no-reference composite for super-resolution evaluation.
    \item \textbf{MUSIQ} \cite{ke_musiq_2021} uses a multi-scale image quality transformer.
          The detail-preserving baseline scores higher than the text-conditioned smoother
          baseline on MUSIQ across all test videos in our preliminary work.
    \item \textbf{DOVER} \cite{wu_dover_2023} combines technical and aesthetic video quality
          dimensions through two disentangled convolutional branches whose outputs are
          combined into a single video-level score. The detail-preserving baseline leads
          the text-conditioned smoother by approximately $+8.75$ DOVER points on the
          per-method mean over the real-SR test set.
\end{itemize}

\subsection{Identified Literature Gap}

Three properties are needed for rigorous long-video SR evaluation but absent from any
single published metric:

\begin{enumerate}
    \item \textbf{No-reference} --- HR ground truth is unavailable in deployment.
    \item \textbf{Multi-time-scale} --- adjacent-frame metrics miss the tOF crossover at
          $k{=}5$--$10$ frames documented in preliminary work.
    \item \textbf{Bias-calibrated} --- three independent learned-representation metrics
          (DINOv2, anomaly ViT, LPIPS) share a structural smoother-output bias; a
          principled aggregation must account for this.
\end{enumerate}

The proposed Long-Range Video Consistency Composite (LR-VCC) metric, described in
Section~\ref{sec:method}, is designed to satisfy all three.
